% Column reference for the GNoME oxide-screen CSVs. Build with:
%   cd docs && pdflatex oxide_screen_columns.tex && pdflatex oxide_screen_columns.tex
\documentclass[11pt,a4paper]{article}

\usepackage[margin=1in]{geometry}
\usepackage{amsmath,amssymb}
\usepackage{booktabs}
\usepackage{xcolor}
\usepackage{tikz}
\usepackage[T1]{fontenc}
\usepackage{lmodern}
\usepackage[colorlinks=true,linkcolor=linknavy,urlcolor=linknavy]{hyperref}

\emergencystretch=3em

\definecolor{linknavy}{HTML}{2A78D6}
\definecolor{inkSoft}{HTML}{52514E}
\definecolor{boxbg}{HTML}{F4F6FA}
\definecolor{boxrule}{HTML}{C9D6EA}

\newcommand{\keybox}[2]{%
  \begin{center}
  \begin{tikzpicture}
    \node[draw=boxrule,fill=boxbg,rounded corners=3pt,line width=0.8pt,
          inner sep=10pt,text width=0.92\textwidth,align=left]
      {\textbf{#1}\\[3pt]\small #2};
  \end{tikzpicture}
  \end{center}}

\title{\vspace{-1.5cm}\textbf{Reading the GNoME Oxide Screen}\\[4pt]
\large A column-by-column reference for the two oxide CSVs}
\author{}
\date{}

\begin{document}
\maketitle
\vspace{-1.3cm}

\noindent
\texttt{15\_filter\_oxides.py} reads the full GNoME screen output ---
\texttt{gnome\_screen\_all.csv}, 33{,}323 candidates scored by
\texttt{13\_screen\_gnome.py} --- and pulls out every row whose formula
contains oxygen, writing two files (both in \texttt{results/cgcnn/}):

\begin{center}
\begin{tabular}{ll}
\toprule
File & Rows \\
\midrule
\texttt{gnome\_oxide\_candidates.csv} & 11{,}612 --- every scored oxide \\
\texttt{gnome\_oxide\_low\_kappa\_candidates.csv} & 1{,}545 --- oxides with $\kappa_L \leq 1$ W~m$^{-1}$~K$^{-1}$ \\
\bottomrule
\end{tabular}
\end{center}

\noindent
Both share the same columns (the low-$\kappa$ file drops
\texttt{is\_low\_kappa\_candidate}, since every row there is one by
definition). This document explains what each column means.

% -----------------------------------------------------------------------
\subsection*{Identity}

\begin{center}
\begin{tabular}{lp{10.3cm}}
\toprule
Column & Meaning \\
\midrule
\texttt{material\_id} & GNoME's own identifier for the material. \\
\texttt{formula} & Reduced chemical formula, e.g.\ \texttt{SbWO4}. \\
\texttt{stoichiometry\_pattern} & A rough structural family, read off the
  \emph{reduced} cation:oxygen ratio --- see the box below. \\
\texttt{is\_low\_kappa\_candidate} & (\texttt{gnome\_oxide\_candidates.csv}
  only) \texttt{True} if this row also appears in the low-$\kappa$ file. \\
\bottomrule
\end{tabular}
\end{center}

\keybox{What \texttt{stoichiometry\_pattern} means}{A pattern-match on the
\emph{ratio} of cations to oxygen after reducing to lowest integers --- not a
verified crystal structure, which would need atomic positions this column
does not look at. \textbf{ABO3-type}: two cations in 1:1 ratio, three oxygens
(perovskite-like). \textbf{AB2O4-type}: cations in 1:2, four oxygens
(spinel-like). \textbf{A2B2O7-type}: cations in 1:1 (reduces from 2:2), seven
oxygens (pyrochlore-like). Anything else with two cations is labelled a
generic ternary oxide; three or more cations become \textit{complex oxide
($n$ elements)}. In this particular screen no exact ABO3 or simple binary
(single-cation) oxide survived the paper's own bandgap/stability/
non-radioactive filters --- GNoME's discovered materials skew toward more
complex, previously-uncatalogued compositions, and simple perovskites are
mostly already well known.}

% -----------------------------------------------------------------------
\subsection*{Structural quantities}

\noindent Read directly from the primitive cell of the parsed structure ---
not predicted, not uncertain.

\begin{center}
\begin{tabular}{lp{10.3cm}}
\toprule
Column & Meaning \\
\midrule
\texttt{Number of Atoms} & Atom count in the \emph{primitive} cell (not
  whatever supercell GNoME's CIF happened to use). \\
\texttt{Volume (A3)} & Primitive-cell volume, \AA$^3$. \\
\texttt{Density (g cm-3)} & Mass $\div$ volume. \\
\texttt{Atomic mass (amu)} & Total composition weight of the primitive cell. \\
\bottomrule
\end{tabular}
\end{center}

% -----------------------------------------------------------------------
\subsection*{Model predictions}

\noindent The actual machine-learning output --- the mean of the 3-model
CGCNN ensemble, in linear GPa (converted back from the $\log_{10}$ space the
network is trained in).

\begin{center}
\begin{tabular}{lp{10.3cm}}
\toprule
Column & Meaning \\
\midrule
\texttt{K\_VRH\_pred} & Predicted bulk modulus $K$ (GPa). \\
\texttt{G\_VRH\_pred} & Predicted shear modulus $G$ (GPa). \\
\texttt{K\_VRH\_spread\_log10} & How much the 3 ensemble members
  \emph{disagree} on $K$, in $\log_{10}$(GPa) --- a per-crystal uncertainty
  estimate, not a physical property. \\
\texttt{G\_VRH\_spread\_log10} & Same disagreement measure for $G$. \\
\bottomrule
\end{tabular}
\end{center}

% -----------------------------------------------------------------------
\subsection*{Derived physics (the Slack model)}

\noindent Everything below is computed from $K$ and $G$ above via the
paper's unchanged Slack-model formulas
(\texttt{scripts/cgcnn/07\_predict\_kappa.py}, \texttt{slack\_physics()}).

\begin{center}
\begin{tabular}{lp{10.3cm}}
\toprule
Column & Meaning \\
\midrule
\texttt{Poisson ratio} & From the ratio of longitudinal to transverse sound
  velocity, $v_{\mathrm{long}} = \sqrt{(K + \tfrac{4}{3}G)/\rho}$,
  $v_{\mathrm{trans}} = \sqrt{G/\rho}$. \\
\texttt{Gruneisen parameter} & How anharmonic the lattice vibrations are,
  derived from the Poisson ratio. Larger $\Rightarrow$ phonons scatter off
  each other more strongly $\Rightarrow$ generally lower $\kappa_L$. \\
\texttt{Kappa\_cal (W m-1 K-1)} & \textbf{The predicted lattice thermal
  conductivity} --- the number this whole screen thresholds on. \\
\bottomrule
\end{tabular}
\end{center}

\begin{center}
$\kappa_L = \dfrac{G \cdot v_{\mathrm{sound}} \cdot V^{1/3}}
{N_{\mathrm{atoms}} \cdot 300} \cdot e^{-\gamma}$
\quad (Eq.\ 2 of the PINK paper; $\delta=1$)
\end{center}

% -----------------------------------------------------------------------
\subsection*{Uncertainty interval (Monte Carlo)}

\begin{center}
\begin{tabular}{lp{10.3cm}}
\toprule
Column & Meaning \\
\midrule
\texttt{Kappa\_cal\_p05} & 5th percentile of $\kappa_L$ across 2{,}000 Monte
  Carlo draws. \\
\texttt{Kappa\_cal\_p50} & Median of the same 2{,}000 draws --- sits close to
  \texttt{Kappa\_cal} above. \\
\texttt{Kappa\_cal\_p95} & 95th percentile. \\
\bottomrule
\end{tabular}
\end{center}

\keybox{How the interval is built, and how to read it}{$K$ and $G$ are each
redrawn 2{,}000 times from a normal distribution centred on the ensemble
prediction, with width equal to \texttt{K\_VRH\_spread\_log10} /
\texttt{G\_VRH\_spread\_log10} --- i.e.\ how much the 3 ensemble members
disagreed. Structural quantities stay fixed; only the \emph{predicted} moduli
are uncertain. Every draw goes through the identical Slack-model formula used
for the point estimate, so p05/p50/p95 are just percentiles of the resulting
2{,}000 $\kappa_L$ values --- consistent with \texttt{Kappa\_cal} by
construction, not a separate approximation. \textbf{A tight p05--p95 range
means the ensemble agreed and the prediction is trustworthy; a wide range
means the 3 models disagreed and the candidate is worth a real DFT check
before believing it.} This is something a single point-estimate model (like
the paper's own released checkpoint) cannot report at all.}

\end{document}
